\section{Related Work}

General embedding and multilingual retrieval benchmarks, including MTEB, MMTEB,
MIRACL, NeuCLIR, and CLIRMatrix, enable systematic comparison of retrieval and
embedding models across languages
\citep{miracl2023,neuclir2023,muennighoff2023mteb,mmteb2025,clirmatrix2020}.
Chemistry-specific resources such as ChemTEB, ChemLit-QA, ChemComp,
ChemKGMultiHopQA, and ChEmbed extend evaluation to chemical embeddings, RAG, and
scientific reasoning
\citep{chemteb2024,chemlitqa2024,khodadad2026chemcomp,astaraki2026iterativerag,chembed2025}.
However, these lines of work do not jointly provide multilingual chemistry
retrieval data with controlled technical variants, explicit same-language and
cross-language relevance judgments, and chemically plausible hard negatives. This
missing combination motivates the benchmark suite introduced in this paper.

Recent work on cross-lingual retrieval and multilingual RAG shows that retrieval
behavior can depend strongly on language alignment, translation effects, and
same-language preference
\citep{whatdrivesclir2025,crosslingualcost2025,bordirlines2024,xrag2025,nepotism2025}.
These findings expose a gap for multilingual chemistry retrieval: aggregate
recall alone does not show whether a model retrieves the right evidence across
languages or mainly benefits from easier same-language matches. Patent
publication variants are useful for studying this gap because they provide
comparable multilingual technical evidence under tighter content control than
independently collected documents across languages.

Patent retrieval has a long evaluation history, including NTCIR and CLEF-IP
\citep{prime2002,clefip2013}, while recent work has revisited patent
representation learning and embedding evaluation
\citep{paecter2024,patentembeddings2026}. Because patent publications often
contain related technical disclosures across languages, they can serve as a
source of comparable multilingual technical evidence beyond prior-art search
alone.

Chemical ontology and alias resources, including ChEBI \citep{chebi2016}, also
make it possible to define hard negatives that are plausible in chemistry search
but incorrect for a target query. Such chemistry-aware retrieval tests complement
ontology-grounded extraction work such as CEAR \citep{cear2024}: CEAR extracts
chemical entities and roles, whereas retrieval evaluation asks whether models
distinguish intended evidence from chemically related but incorrect documents.
